\subsection{比较关系的求导}

\pandocbounded{\includegraphics[keepaspectratio]{images/part002_390c97e1.jpg}}

命题 1 (V, p.~228) 与命题 6 (V, p.~230) 没有逆命题：两个函数之间的比较关系 \(\mathbf{f} \preccurlyeq g\) 、 \(\mathbf{f} \ll g\) 、 \(\mathbf{f} \sim \mathbf{c}g\) 的存在（这两个函数在 \(+\infty\) 的一个邻域上可导）并不蕴含它们的导数之间有同样的比较关系，即使对实单调函数 \(f\) 与 \(g\) 亦然。

例如，函数 \(x^{2} + x\sin x + \cos x\) 是单调的且等价于 \(x^{2}\) ，但它的导数 \(x(2 + \cos x)\) 不等价于 \(2x\) 。

另一方面，当先验地假设所考虑函数的导数是可比较的 (V, p.~217) 时，可以导出比较关系。一般地，我们将说定义在区间 \([a, +\infty]\) 上的两个实函数 f 与 g 在 \(+\infty\) 的一个邻域上是 k 阶可比较的，如果在 \(+\infty\) 的一个邻域上，它们各自除在一个可数集的点处外都有规则的 \(k^{th}\) 导数，并且在这个邻域上 \(f^{(k)}\) 与 \(g^{(k)}\) 具有常号（在其有定义的集上）且是可比较的。

我们约定称两个可比较的实函数（V, p.~217）为 0 阶可比较的。

命题 7. 若两个实函数 \(f, g\) 是 1 阶可比较的，则它们是可比较的；此外，关系 \(f \ll g\) （相应地， \(f \sim cg\) ， \(c\) 为常数）蕴含 \(f' \ll g'\) （相应地， \(f' \sim cg'\) ），除非 \(g\) 等价于一个非零常数。

现在，由于 \(f'\) 与 \(g'\) 在一个区间 \([x_0, +\infty[\) 上具有常号， \(f\) 与 \(g\) 都在这个区间上单调，故当 \(x\) 趋于 \(+\infty\) 时趋于一个有限或无穷的极限。显然， \(f\) 与 \(g\) 当 \(x\) 趋于 \(+\infty\) 时是可比较的（若这两个极限之一有限且 \(\neq 0\) ，或者一个为零而另一个为无穷）。若 \(f\) 与 \(g\) 都趋于 0，可以写 \(f(x) = -\int_x^{+\infty} f'(t) dt\) ， \(g(x) = -\int_x^{+\infty} g'(t) dt\) ；由于 \(f'\) 与 \(g'\) 可比较， \(f\) 与 \(g\) 亦然，且 \(f\) 与 \(g\) 之间的比较关系等同于 \(f'\) 与 \(g'\) 之间的比较关系（由 V, p.~20 的命题 6）。类似地，若 \(f\) 与 \(g\) 都有一个无穷极限，则有 \(f(x) = f(x_0) + \int_{x_0}^x f'(t) dt\) ， \(g(x) = g(x_0) + \int_{x_0}^x g'(t) dt\) ；V, p.~230 的命题 6 仍然表明 \(f\) 与 \(g\) 可比较，且 \(f\) 与 \(g\) 之间的比较关系等同于 \(f'\) 与 \(g'\) 之间的比较关系。为完成证明，剩下考虑 \(g\) 趋于 \(\pm \infty\) 而 \(f\) 趋于一个常数的情形；这时不可能有 \(f' \succcurlyeq g'\) ，因为那样由 V, p.~218 的命题 1 将可推出积分 \(\int_{x_0}^{+\infty} g'(t) dt\) 收敛；由于已假设 \(f'\) 与 \(g'\) 可比较，必有 \(f' \ll g'\) 。

推论. 若两个实函数 \(f, g\) 是 \(k \geqslant 1\) 阶可比较的，则它们是 \(p\) 阶可比较的（对 \(0 \leqslant p \leqslant k\) ）；此外，关系 \(f \ll g\) （相应地， \(f \sim cg\) ）蕴含 \(f^{(k)} \ll g^{(k)}\) （相应地， \(f^{(k)} \sim cg^{(k)}\) ），除非诸导数 \(g^{(p)}\) 之一（ \(0 \leqslant p \leqslant k - 1\) ）等价于一个常数 \(\neq 0\) 。

事实上，由于 \(f^{(k)}\) 与 \(g^{(k)}\) 在一个区间 \([x_0, +\infty[\) 上具有常号，由此得出 \(f^{(k-1)}\) 与 \(g^{(k-1)}\) 在这个区间上单调，故在 \(+\infty\) 的一个邻域上具有常号；此外，V, p.~232 的命题 7 表明 \(f^{(k-1)}\) 与 \(g^{(k-1)}\) 可比较，于是推论由递归地应用命题 7 而得。

注. 1) 命题 7 中对 \(g\) 的限制是本质的。例如，有 \(\frac{1}{x} \ll 1 + \frac{1}{x}\) ，尽管两边的导数是等价的；类似地 \(1 + \frac{1}{x} \sim 1 + \frac{1}{x^2}\) ，但 \(1 / x^2 \gg 2 / x^3\) 。

\begin{enumerate}
\def\labelenumi{\arabic{enumi})}
\setcounter{enumi}{1}
\item
  尽管 \(f\) 与 \(g\) 是 \(k\) 阶可比较的，一个函数 \(f_{1}\) （等价于 \(f\) 的）不必是 \(k\) 阶可比较的（与 \(g\) ）；但若假设 \(f_{1}\) 是 \(k\) 阶可比较的（与 \(f\) ），且诸导数 \(f^{(p)}\) （ \(0 \leqslant p \leqslant k - 1\) ）无一等价于一个非零常数，则它将是如此。
\item
  若 f 与 g 是 k 阶可比较的，这对 hf 与 hg 未必成立，即使对像 \(h(x) = x\) 这样简单的单调函数 h 亦然 (V, p.~260, exerc. 3)；同样，1/f 与 1/g 也不必然是 k 阶可比较的 (V, p.~259, exerc. 1)。
\end{enumerate}

